\chapter{L'oscillateur harmonique quantique}\label{ch:b3:harmonic-oscillator}

Presque rien dans la nature n'est exactement un oscillateur
harmonique, et presque tout en est un approximativement~: tout système
écarté d'un équilibre stable --- la liaison d'une molécule, un atome
dans un cristal, un pont, un mode du champ électromagnétique ---
ressent une force de rappel proportionnelle à l'écart, parce que tout
potentiel régulier est une parabole au fond de son puits. Qui résout
une fois l'oscillateur quantique résout donc les petites vibrations du
monde entier. Ce chapitre le résout dans le style algébrique devenu la
signature de la mécanique quantique~: deux \omterm{def:b3:harmonic-oscillator:ladder}{opérateurs d'échelle}
montent et descendent un escalier parfaitement régulier de niveaux $\big(n +
\tfrac12\big)\hbar\omega$, la demi-marche du bas --- l'\omterm{thm:b3:harmonic-oscillator:spectrum}{énergie de
point zéro} --- étant un théorème et non une option. Les conséquences
vont de la raison pour laquelle l'hélium ne gèle jamais à la façon dont
une molécule de dioxyde de carbone, sonnant à son propre $\hbar\omega$,
intercepte la chaleur qui quitte la Terre.

\section{L'échelle}

\begin{definition}[Opérateurs d'échelle]\label{def:b3:harmonic-oscillator:ladder}
\index{opérateurs d'échelle}\index{opérateur d'annihilation}\index{opérateur de création}
Pour $\hat H = \hat p^2/2m + \tfrac12 m\omega^2\hat x^2$, introduisons
le couple sans dimension et non hermitien
\[
\hat a = \sqrt{\frac{m\omega}{2\hbar}}\Big(\hat x +
\frac{\iu\hat p}{m\omega}\Big) , \qquad
\hat a^\dagger = \sqrt{\frac{m\omega}{2\hbar}}\Big(\hat x -
\frac{\iu\hat p}{m\omega}\Big) .
\]
À partir de $[\hat x, \hat p] = \iu\hbar$~:
\[
[\hat a, \hat a^\dagger] = 1 , \qquad
\hat H = \hbar\omega\Big(\hat N + \tfrac12\Big) , \quad
\hat N = \hat a^\dagger\hat a .
\]
$\hat N$ est hermitien~; ses \omterm{def:b3:quantum-formalism:observable}{valeurs propres} compteront des
\emph{quanta}, et $\hat a$, $\hat a^\dagger$ --- les opérateurs
d'\emph{annihilation} et de \emph{création} --- en retireront et en
ajouteront un.
\end{definition}

\begin{theorem}[Le spectre, par la seule algèbre]\label{thm:b3:harmonic-oscillator:spectrum}
\index{énergie de point zéro}
Les \omterm{def:b3:quantum-formalism:observable}{valeurs propres} de $\hat H$ sont exactement
\[
E_n = \Big(n + \tfrac12\Big)\hbar\omega , \qquad n = 0, 1, 2, \dots
\]
--- une échelle infinie de marches égales $\hbar\omega$ au-dessus d'un
niveau fondamental qui n'est \emph{pas} nul~: l'\emph{\omterm{thm:b3:harmonic-oscillator:spectrum}{énergie de point
zéro}} $\tfrac12\hbar\omega$. Les états propres normés $\ket n$ sont
reliés par
\[
\hat a\ket n = \sqrt n\,\ket{n - 1} , \qquad
\hat a^\dagger\ket n = \sqrt{n + 1}\,\ket{n + 1} ,
\qquad
\ket n = \frac{(\hat a^\dagger)^n}{\sqrt{n!}}\,\ket0 .
\]
\end{theorem}

\begin{proof}
D'après le \omterm{def:b3:quantum-formalism:commutator}{commutateur}, $\hat N\hat a = \hat a(\hat N - 1)$~: si
$\ket\nu$ a pour $\hat N$ la \omterm{def:b3:quantum-formalism:observable}{valeur propre} $\nu$, alors $\hat a\ket\nu$
est un vecteur propre de valeur $\nu - 1$ (ou le vecteur nul). Chaque
descente n'est permise que tant que $\nu \ge 0$, puisque $\nu = \braket{\nu}{\hat
N\nu} = \|\hat a\ket\nu\|^2 \ge 0$~; la descente doit donc
s'interrompre, et elle ne s'interrompt que sur un état vérifiant $\hat a\ket{\nu_0}
= 0$, d'où $\nu_0 = 0$. Ainsi $\nu$ parcourt les entiers positifs ou
nuls. Les normalisations découlent de $\|\hat a\ket n\|^2 = n$ et de
$\|\hat a^\dagger\ket n\|^2 = n + 1$.
\end{proof}

\begin{omfigure}
\begin{tikzpicture}[font=\small, >=Stealth]
  \draw[->] (0,-0.3) -- (0,5.4);
  \node[anchor=south] at (0,5.4) {$E$};
  \foreach \n in {0,1,2,3,4}
    {\draw[very thick, omDef] (0.5,{0.55+\n*1.05}) -- (3.3,{0.55+\n*1.05});
     \node[anchor=west] at (3.45,{0.55+\n*1.05})
       {$\ket{\n}$\; $E = (\n + \tfrac12)\hbar\omega$};}
  \draw[->, omThm, very thick] (1.35,0.62) -- (1.35,1.53);
  \node[omThm, anchor=east] at (1.28,1.08) {$\hat a^\dagger$};
  \draw[->, omExo, very thick] (2.45,1.53) -- (2.45,0.62);
  \node[omExo, anchor=west] at (2.52,1.08) {$\hat a$};
  \draw[<->] (0.15,0) -- (0.15,0.55);
  \node[anchor=west, font=\footnotesize] at (0.28,0.22) {$\tfrac12\hbar\omega$};
  \draw[dashed] (0,0) -- (3.3,0);
\end{tikzpicture}

{\small L'échelle de l'oscillateur~: des marches égales $\hbar\omega$,
gravies par $\hat a^\dagger$ et descendues par $\hat a$, reposant sur
un plancher situé une demi-marche au-dessus de l'\omterm{def:b3:relativistic-dynamics:momentum-energy}{énergie de masse}
classique.}
\end{omfigure}

\begin{proposition}[Les états dans l'espace]\label{prop:b3:harmonic-oscillator:wavefunctions}
L'état fondamental est la gaussienne
\[
\varphi_0(x) = \Big(\frac{m\omega}{\pi\hbar}\Big)^{1/4}
\eu^{-m\omega x^2/2\hbar} ,
\]
obtenue en résolvant l'équation du \emph{premier ordre} $\hat a\varphi_0
= 0$~; elle sature l'inégalité de Heisenberg, $\Delta x\,\Delta p =
\hbar/2$, avec $\Delta x = \sqrt{\hbar/2m\omega}$. L'application
répétée de $\hat a^\dagger$ engendre $\varphi_n$~: un polynôme de degré
$n$ (à $n$ nœuds) multiplié par la même gaussienne. Pour $n$ grand,
$|\varphi_n|^2$ oscille rapidement autour de la distribution classique
des temps de séjour, maximale près des points de rebroussement --- le
principe de correspondance en image.
\end{proposition}

\begin{proof}[Démonstration partielle]
$\hat a\varphi_0 = 0$ s'écrit $\varphi_0' = -(m\omega/\hbar)x
\varphi_0$~: la gaussienne, normée. Saturation~: la gaussienne est le
cas d'égalité de \cref{thm:b3:quantum-formalism:uncertainty}. La
structure polynomiale découle de ce que $\hat a^\dagger$ est du premier
ordre en $x$ et $\dd/\dd x$~; l'énoncé sur les grands $n$ est vérifié
dans \cref{exo:b3:harmonic-oscillator:12}.
\end{proof}

\begin{omfigure}
\begin{tikzpicture}
  \begin{axis}[omaxis, axis x line=bottom, axis y line=left,
    width=6.6cm, height=5.2cm, xmin=-4, xmax=4, ymin=0, ymax=3.6,
    xlabel={$x$ (en unités de $\sqrt{\hbar/m\omega}$)}, ylabel={},
    xlabel style={at={(axis description cs:0.5,-0.1)}, anchor=north},
    xtick={-2,0,2}, ytick=\empty, clip=false]
    \addplot[omDef, thick, domain=-4:4, samples=90]
      {0.56*exp(-x*x)+0.25};
    \addplot[omThm, thick, domain=-4:4, samples=90]
      {1.13*x*x*exp(-x*x)+1.3};
    \addplot[omExo, thick, domain=-4:4, samples=120]
      {0.28*(2*x*x-1)^2*exp(-x*x)+2.35};
    \node[omDef, font=\scriptsize, anchor=west] at (axis cs:2.4,0.5) {$|\varphi_0|^2$};
    \node[omThm, font=\scriptsize, anchor=west] at (axis cs:2.4,1.55) {$|\varphi_1|^2$};
    \node[omExo, font=\scriptsize, anchor=west] at (axis cs:2.4,2.6) {$|\varphi_2|^2$};
  \end{axis}
  \begin{axis}[omaxis, axis x line=bottom, axis y line=left,
    width=6.6cm, height=5.2cm, xmin=-5.2, xmax=5.2, ymin=0, ymax=0.45,
    xshift=7.1cm,
    xlabel={$x$}, ylabel={},
    xlabel style={at={(axis description cs:0.5,-0.1)}, anchor=north},
    xtick=\empty, ytick=\empty, clip=false]
    \addplot[omDef, domain=-4.6:4.6, samples=300]
      {(64*x^6-480*x^4+720*x^2-120)^2*exp(-x*x)/81680};
    \addplot[omThm, very thick, dashed, domain=-3.54:3.54, samples=120]
      {0.3183/sqrt(max(0.35,13-x*x))};
    \node[omThm, font=\scriptsize, anchor=south] at (axis cs:0,0.36) {temps de séjour classique};
    \node[omDef, font=\scriptsize, anchor=west] at (axis cs:1.9,0.42) {$|\varphi_6|^2$};
  \end{axis}
\end{tikzpicture}

{\small À gauche~: les densités de probabilité les plus basses
(décalées verticalement)~: une gaussienne sans nœud, puis $n$ nœuds
pour $\varphi_n$. À droite~: aux grands $n$ la densité quantique
oscille autour de la distribution classique, qui s'accumule aux points
de rebroussement où la masse oscillante s'attarde.}
\end{omfigure}

\section{L'énergie de point zéro est réelle}

\begin{remark}[Pourquoi le plancher ne peut pas être plus bas]\label{rem:b3:harmonic-oscillator:zeropoint}
Un état d'énergie nulle exigerait $\langle\hat p^2\rangle =
\langle\hat x^2\rangle = 0$~: parfaitement immobile \emph{et}
parfaitement centré, ce qu'interdit $\Delta x\,\Delta p \ge \hbar/2$.
L'état fondamental est le meilleur compromis que l'inégalité autorise,
et $\tfrac12\hbar\omega$ en est le loyer. Ce n'est pas une constante de
comptabilité~: le mouvement de point zéro brouille les diagrammes de
diffraction des rayons X au zéro absolu~; il donne aux isotopes plus
légers des liaisons effectives plus faibles (H$_2$ et D$_2$ se
dissocient à des énergies mesurablement différentes depuis le même
puits électronique)~; et dans l'hélium il est si violent --- atomes
légers, attraction faible --- que le liquide ne gèle \emph{jamais} sous
sa propre pression de vapeur~: le seul élément encore liquide au zéro
absolu, qui ne se solidifie qu'avec l'aide de 25 atmosphères.
\end{remark}

\begin{example}[Échelles de $\hbar\omega$]\label{ex:b3:harmonic-oscillator:scales}
Liaison du monoxyde de carbone ($\omega = \qty{4.1e14}{rad/s}$)~: $\hbar\omega =
\qty{0.27}{eV}$, dix fois le $k_{\text{B}}T$ ambiant --- les vibrations
moléculaires sont gelées dans l'air de tous les jours, ce qui explique
que les capacités thermiques des gaz diatomiques aient intrigué le
dix-neuvième siècle. Un pendule ($\omega = \qty{5}{rad/s}$)~:
$\hbar\omega = \qty{3e-15}{eV}$, désespérément hors de portée --- la
limite de correspondance. Un miroir de LIGO de \qty{40}{kg}, suspendu
de façon à osciller à \qty{100}{Hz}, a comme mode d'oscillateur l'amplitude de point zéro $\Delta x =
\sqrt{\hbar/2m\omega} \approx \qty{5e-20}{m}$ --- et l'observatoire
résout couramment des déplacements \emph{à} ce plancher quantique~: le
mouvement de point zéro d'un objet de quarante kilogrammes est
désormais une contrainte d'ingénierie
(\cref{exo:b3:harmonic-oscillator:3}).
\end{example}

\section{États cohérents~: le visage classique de l'oscillateur quantique}

\begin{definition}[États cohérents]\label{def:b3:harmonic-oscillator:coherent}
\index{état cohérent}
Un \emph{état cohérent} $\ket\alpha$ est un vecteur propre de
l'\omterm{def:b3:harmonic-oscillator:ladder}{opérateur d'annihilation}, $\hat a\ket\alpha = \alpha\ket\alpha$, avec
$\alpha$ complexe quelconque. Développé sur l'échelle,
\[
\ket\alpha = \eu^{-|\alpha|^2/2}\sum_{n=0}^\infty
\frac{\alpha^n}{\sqrt{n!}}\,\ket n :
\]
le nombre de quanta suit une loi de Poisson de moyenne $\langle\hat
N\rangle = |\alpha|^2$ et d'écart $\Delta N = |\alpha|$.
\end{definition}

\begin{proposition}[Pourquoi les lasers sont classiques]\label{prop:b3:harmonic-oscillator:coherent-motion}
Un \omterm{def:b3:harmonic-oscillator:coherent}{état cohérent} est la gaussienne de l'état fondamental déplacée dans
l'\omterm{def:b3:hamiltonian-mechanics:hamiltonian}{espace des phases}~; sous l'évolution de l'oscillateur il \emph{reste}
cohérent, $\alpha(t) = \alpha\,\eu^{-\iu\omega t}$~: son centre exécute
exactement le mouvement classique, $\langle\hat x\rangle(t) =
x_0\cos\omega t + \cdots$, tandis que ses largeurs restent à jamais
minimales, $\Delta x\,\Delta p = \hbar/2$ --- un paquet d'ondes qui ne
s'étale jamais. Les \omterm{def:b3:harmonic-oscillator:coherent}{états cohérents} sont la façon dont un oscillateur
quantique se fait passer pour un oscillateur classique~; la lumière
d'un laser est un \omterm{def:b3:harmonic-oscillator:coherent}{état cohérent} d'un mode du champ, son nombre de
photons poissonien (le \emph{bruit de grenaille} de tout
photodétecteur), son champ oscillant comme Maxwell l'a dit.
\end{proposition}

\begin{proof}[Démonstration partielle]
Le développement s'obtient en écrivant $\ket\alpha = \sum c_n\ket n$
dans $\hat a\ket\alpha = \alpha\ket\alpha$~: $c_n = \alpha c_{n-1}/\sqrt
n$. Évolution~: chaque $\ket n$ acquiert $\eu^{-\iu(n + 1/2)\omega t}$,
ce qui se resomme en un \omterm{def:b3:harmonic-oscillator:coherent}{état cohérent} de $\alpha\eu^{-\iu\omega t}$
(à une phase globale près). $\langle\hat x\rangle \propto
\operatorname{Re}\alpha(t)$ d'après $\hat x \propto \hat a +
\hat a^\dagger$. Que l'état soit la gaussienne déplacée est admis ici.
\end{proof}

\begin{omfigure}
\begin{tikzpicture}[font=\small, >=Stealth]
  \draw[->] (-3.1,0) -- (3.4,0) node[right] {$\langle\hat x\rangle$};
  \draw[->] (0,-3.0) -- (0,3.2) node[above] {$\langle\hat p\rangle$};
  \draw[gray, dashed] (0,0) circle (2.2);
  \fill[omDef!25, draw=omDef, thick] (2.2,0) circle (0.42);
  \fill[omDef!25, draw=omDef, thick] (50:2.2) circle (0.42);
  \fill[omDef!25, draw=omDef, thick] (100:2.2) circle (0.42);
  \draw[->, thick] (2.53,0.9) arc (20:40:2.6);
  \node[anchor=west] at (2.62,1.35) {$\omega$};
  \fill[omThm!30, draw=omThm, thick] (0,0) circle (0.42);
  \node[omThm, font=\footnotesize, anchor=east] at (-0.3,-0.68) {$\ket0$};
  \node[omDef, font=\footnotesize, anchor=west] at (2.3,-0.62) {$\ket\alpha$};
  \node[font=\footnotesize, anchor=west] at (-3.05,2.75)
    {tache minimale $\Delta x\,\Delta p = \hbar/2$, tournant classiquement};
\end{tikzpicture}

{\small Portrait dans l'\omterm{def:b3:hamiltonian-mechanics:hamiltonian}{espace des phases} (à comparer au
\cref{ch:b3:hamiltonian-mechanics})~: l'état fondamental est une tache
d'indétermination minimale à l'origine~; un \omterm{def:b3:harmonic-oscillator:coherent}{état cohérent} est la même
tache déplacée, tournant à $\omega$ sans se déformer --- la meilleure
imitation d'une oscillation classique que la mécanique quantique
sache produire.}
\end{omfigure}

\begin{method}[Algèbre de l'oscillateur]\label{met:b3:harmonic-oscillator:method}
(1) Exprimer ce qui est demandé en $\hat a$, $\hat a^\dagger$~: $\hat x
= \sqrt{\hbar/2m\omega}\,(\hat a + \hat a^\dagger)$, $\hat p =
\iu\sqrt{m\hbar\omega/2}\,(\hat a^\dagger - \hat a)$. (2) Faire passer
les $\hat a$ à droite avec $[\hat a, \hat a^\dagger] = 1$~; utiliser
$\hat a\ket0 = 0$. (3) Éléments de matrice~: $\hat x$ ne relie que des
barreaux voisins --- la règle de sélection $\Delta n = \pm1$ de la
spectroscopie vibrationnelle. (4) Tout \omterm{def:b3:hamiltonian-mechanics:hamiltonian}{hamiltonien} quadratique
(oscillateurs couplés, circuits, modes du champ) se diagonalise en
échelles indépendantes~: trouver d'abord les \omterm{prop:b3:lagrangian-mechanics:modes}{modes propres}, puis
quantifier chacun. (5) Les nombres d'abord~: $\hbar\omega$ comparé à
$k_{\text{B}}T$ décide du caractère quantique ou classique du système
avant toute algèbre.
\end{method}

\section{Exercices}

\begin{exercise}[$\star$]\label{exo:b3:harmonic-oscillator:1}
(a) Vérifier $[\hat a, \hat a^\dagger] = 1$ à partir de $[\hat x, \hat p] =
\iu\hbar$. (b) Vérifier $\hat H = \hbar\omega(\hat a^\dagger\hat a +
\tfrac12)$ par développement direct. (c) Inverser pour exprimer
$\hat x$ et $\hat p$. (d) Montrer que $\hat N$ est hermitien et
expliquer pourquoi $\hat a$ seul ne saurait être une observable.
\end{exercise}

\begin{exercise}[$\star$]\label{exo:b3:harmonic-oscillator:2}
Sur l'état propre $\ket n$~: (a) montrer que $\langle\hat x\rangle =
\langle\hat p\rangle = 0$~; (b) calculer $\langle\hat x^2\rangle$ et
$\langle\hat p^2\rangle$~; (c) en déduire $\Delta x\,\Delta p = (n +
\tfrac12)\hbar$~; (d) vérifier le rapport du viriel $\langle E_k\rangle =
\langle E_p\rangle$.
\end{exercise}

\begin{exercise}[$\star$]\label{exo:b3:harmonic-oscillator:3}
Amplitudes de point zéro $\Delta x = \sqrt{\hbar/2m\omega}$~: les calculer
pour (a) la molécule CO ($\mu = \qty{1.14e-26}{kg}$, $\omega =
\qty{4.1e14}{rad/s}$), en comparaison de la longueur de liaison
\qty{0.11}{nm}~; (b) un atome d'hydrogène dans un solide ($m =
\qty{1.7e-27}{kg}$, $\hbar\omega = \qty{0.1}{eV}$)~; (c) un miroir de
LIGO ($m = \qty{40}{kg}$, $f = \qty{100}{Hz}$)~; (d) classer les trois
en fractions de la taille de leur système et dites qui est
«~quantique~».
\end{exercise}

\begin{exercise}[$\star$]\label{exo:b3:harmonic-oscillator:4}
(a) À l'aide des relations d'échelle, calculer $\bra m\hat x\ket n$ et
montrer qu'il s'annule sauf si $m = n \pm 1$. (b) En déduire la règle
de sélection vibrationnelle $\Delta n = \pm1$ pour l'absorption de
lumière (le couplage est $\propto\hat x$). (c) Pourquoi une molécule
hétéronucléaire (CO) absorbe-t-elle l'infrarouge alors que N$_2$ ne le
fait pas~? (d) Les molécules réelles montrent de faibles raies
«~harmoniques~» vers $\approx 2\hbar\omega$~: que révèle cela sur le
potentiel~?
\end{exercise}

\begin{exercise}[$\star\star$]\label{exo:b3:harmonic-oscillator:5}
(a) Résoudre $\hat a\varphi_0 = 0$ comme équation différentielle et
normaliser. (b) Engendrer $\varphi_1$ et $\varphi_2$ en appliquant
$\hat a^\dagger$. (c) Vérifier $\varphi_1 \perp \varphi_0$ par la seule
parité. (d) Tracer les trois densités et vérifier le nombre de nœuds.
\end{exercise}

\begin{exercise}[$\star\star$]\label{exo:b3:harmonic-oscillator:6}
Un avant-goût de Boltzmann. Un ensemble d'oscillateurs identiques à la
température $T$ occupe le niveau $n$ avec la probabilité $\propto
\eu^{-E_n/k_{\text{B}}T}$ (volume de deuxième année, facteur de
Boltzmann). (a) Montrer que le nombre quantique moyen vaut $\langle n\rangle =
1/(\eu^{\hbar\omega/k_{\text{B}}T} - 1)$. (b) L'évaluer pour la
vibration de CO à \qty{300}{K} puis à \qty{2000}{K}. (c) Montrer que
l'énergie moyenne tend vers $k_{\text{B}}T$ à haute température
(équipartition retrouvée) et vers $\tfrac12\hbar\omega$ à basse. (d) À
quelle température une vibration à \qty{15}{\micro m} (la flexion du
dioxyde de carbone) porte-t-elle $\langle n\rangle = 0.1$~?
\end{exercise}

\begin{exercise}[$\star\star$]\label{exo:b3:harmonic-oscillator:7}
Statistique des \omterm{def:b3:harmonic-oscillator:coherent}{états cohérents}. (a) À partir du développement de
$\ket\alpha$, montrer que $\mathcal P(n)$ est poissonien de moyenne
$|\alpha|^2$. (b) Montrer que $\langle\hat N\rangle = |\alpha|^2$ et
$\Delta N = |\alpha|$. (c) Un laser de \qty{1}{mW} à \qty{633}{nm}~:
photons par seconde, et fluctuation relative
$\Delta N/\langle N\rangle$ en une seconde. (d) Bruit de grenaille~:
montrer que le rapport bruit sur signal du photocourant décroît comme
$1/\sqrt{\langle N\rangle}$ --- pourquoi les faisceaux intenses
paraissent lisses.
\end{exercise}

\begin{exercise}[$\star\star$]\label{exo:b3:harmonic-oscillator:8}
Évolution d'un \omterm{def:b3:harmonic-oscillator:coherent}{état cohérent}. (a) Appliquer les phases d'évolution au
développement et montrer que l'on obtient $\ket{\alpha(t)}$ avec $\alpha(t) =
\alpha\eu^{-\iu\omega t}$ (à une phase globale près). (b) En déduire
$\langle\hat x\rangle(t)$ et $\langle\hat p\rangle(t)$ et comparer à la
solution classique. (c) Pourquoi un état propre de l'énergie, quoique
stationnaire, ne décrit-il \emph{pas} un pendule qui oscille --- quelle
propriété de l'\omterm{def:b3:harmonic-oscillator:coherent}{état cohérent} le fait~? (d) Estimer $|\alpha|$ pour un
pendule réel (\qty{10}{g}, \qty{10}{cm} d'amplitude, \qty{1}{Hz}) et
commenter.
\end{exercise}

\begin{exercise}[$\star\star$]\label{exo:b3:harmonic-oscillator:9}
Le circuit LC quantique. Une boucle supraconductrice de $L =
\qty{10}{nH}$ et $C = \qty{0.4}{pF}$ fait osciller charge et flux
comme $x$ et $p$. (a) Sa fréquence de résonance (volume de première
année) et le quantum $\hbar\omega$ en $\unit{\micro eV}$. (b)
Au-dessous de quelle température a-t-on $k_{\text{B}}T \ll \hbar\omega$,
de sorte que le circuit siège dans son état fondamental~? (c) Pourquoi
les qubits supraconducteurs fonctionnent-ils dans des réfrigérateurs à
dilution à \qty{20}{mK}~? (d) La fluctuation de tension de point zéro
$\Delta V = \Delta q/C$ avec $\Delta q =
\sqrt{\hbar\omega C/2}$~: l'évaluer.
\end{exercise}

\begin{exercise}[$\star\star\star$]\label{exo:b3:harmonic-oscillator:10}
Van der Waals à partir du mouvement de point zéro. Modélisons deux
atomes neutres distants de $R$ par deux oscillateurs dipolaires
identiques (charge $e$, masse $m$, pulsation $\omega_0$) dont les
écarts sont couplés par l'énergie dipolaire $\lambda\,x_1x_2$ avec
$\lambda = e^2/2\pi\varepsilon_0R^3$. (a) Montrer que les coordonnées
normales $(x_1 \pm x_2)/\sqrt2$ oscillent à
$\omega_\pm = \omega_0\sqrt{1 \pm \lambda/m\omega_0^2}$. (b) L'énergie
fondamentale vaut $\tfrac\hbar2(\omega_+ + \omega_-)$~: développer au
second ordre en $\lambda$ et montrer que l'énergie d'interaction est
\[
\Delta E = -\frac{\hbar\lambda^2}{8m^2\omega_0^3}
\ \propto\ -\frac{1}{R^6} .
\]
(c) Pourquoi cette attraction est-elle universelle --- présente entre
des atomes dépourvus de tout dipôle permanent~? (d) La loi en $1/R^6$
est la force de van der Waals des corrections de gaz réel du volume de
première année~: qu'est-ce qui, microscopiquement, «~fluctue~» --- et
qu'adviendrait-il de cette force dans un monde où $\hbar = 0$~?
\end{exercise}

\begin{exercise}[$\star\star\star$]\label{exo:b3:harmonic-oscillator:11}
Bandes latérales d'un ion piégé. Un ion dans un piège harmonique ($f =
\qty{1}{MHz}$) absorbe la lumière d'un laser. Comme l'ion bouge, son
spectre d'absorption montre la raie électronique à $\nu_0$ flanquée de
raies à $\nu_0 \pm f$~: des transitions qui changent le nombre
quantique \emph{de mouvement} de $\mp1$ en même temps que le nombre
électronique. (a) Que vaut $\hbar\omega_{\text{piège}}$ en
$\unit{neV}$, et pourquoi résoudre les bandes latérales demande-t-il
une raie très fine~? (b) Exciter la raie $\nu_0 - f$ retire un quantum
de mouvement par cycle~: expliquer le \emph{refroidissement par bande
latérale} jusqu'à l'état fondamental. (c) Une fois $\langle n\rangle \approx
0$, la bande latérale basse disparaît entièrement --- pourquoi cette
asymétrie prouve-t-elle que l'état fondamental quantique est atteint~?
(d) C'est ainsi que l'état fondamental de mouvement d'un atome unique
--- et, par d'autres moyens, celui des miroirs de LIGO à l'échelle du
kilogramme --- est certifié~: dites quelle «~température~» l'ion a
atteinte pour $f = \qty{1}{MHz}$ et $\langle n\rangle = 0.05$.
\end{exercise}

\begin{exercise}[$\star\star\star$]\label{exo:b3:harmonic-oscillator:12}
La limite classique, quantitativement. Un oscillateur classique
d'amplitude $A$ passe dans $[x, x + \dd x]$ la fraction $\dd
t/T = \dd x/\pi\sqrt{A^2 - x^2}$. (a) Obtenir cette distribution des
temps de séjour. (b) Pour l'état quantique $n$, tirer $A_n$ de $E_n =
\tfrac12 m\omega^2A_n^2$ et comparer la distribution classique au
$|\varphi_n|^2$ moyenné localement (qui est donné)~: où s'accordent-ils
et où doivent-ils différer~? (c) Montrer que l'écart relatif entre
niveaux voisins, $\Delta E/E$, s'annule comme $1/n$~: l'énergie devient
effectivement continue. (d) Pour le pendule de
\cref{exo:b3:harmonic-oscillator:8}(d), estimer $n$ et $\Delta
E/E$, et conclure l'argument de correspondance en une phrase.
\end{exercise}

\section{Problème~: la molécule qui réchauffe la Terre}

\begin{problem}[{Problème du week-end --- l'échelle quantique du
dioxyde de carbone et l'effet de serre}]\label{pb:b3:harmonic-oscillator:1}
Une molécule de dioxyde de carbone est une chaîne linéaire O=C=O~:
quelques oscillateurs quantifiés. Que le $\hbar\omega$ de son mode de
flexion tombe justement au milieu de la lueur thermique que la Terre
émet vers l'espace explique pourquoi ce gaz-trace --- quatre molécules
sur dix mille --- gouverne le climat de la planète. Données~: nombre
d'onde de flexion $\tilde\nu_2 = \qty{667}{cm^{-1}}$ (l'unité du
spectroscopiste~: $E = hc\tilde\nu$, $\qty{1}{cm^{-1}} =
\qty{1.24e-4}{eV}$)~; élongation antisymétrique $\tilde\nu_3 =
\qty{2349}{cm^{-1}}$~; élongation symétrique $\tilde\nu_1 =
\qty{1388}{cm^{-1}}$~; $k_{\text{B}}T$ à $\qty{288}{K}$ vaut
\qty{24.8}{meV}~; loi de Wien (volume de deuxième année)~: $\lambda_{\max}T =
\qty{2898}{\micro m.K}$.

\medskip\noindent\textbf{Partie I --- Les modes d'une molécule linéaire.}
\begin{enumerate}
  \item Trois atomes ont neuf \omterm{def:b3:lagrangian-mechanics:coordinates}{degrés de liberté}~: combien sont des
        translations de la molécule entière, combien des rotations (la
        molécule est \emph{linéaire}), et combien de vibrations
        restent-ils~?
  \item Décrire les quatre vibrations~: élongation symétrique,
        élongation antisymétrique, et une flexion deux fois dégénérée
        --- pourquoi la flexion vient-elle en double~?
  \item La lumière se couple à un \emph{dipôle électrique} vibrant
        (\cref{exo:b3:harmonic-oscillator:4}). Quels modes de O=C=O
        modulent le moment dipolaire, et lequel est muet dans
        l'infrarouge~?
  \item Convertir les trois nombres d'onde en longueurs d'onde de
        photon, et les situer~: lesquelles sont dans l'infrarouge
        thermique~?
  \item Pourquoi les deux gaz principaux de l'air, N$_2$ et O$_2$,
        n'absorbent-ils pratiquement aucun infrarouge --- et avec
        quelle conséquence pour la transparence de l'atmosphère~?
  \item La vapeur d'eau, coudée et dipolaire, absorbe sur une grande
        part de l'infrarouge~: dans quelle «~fenêtre~» spectrale la
        flexion du dioxyde de carbone opère-t-elle largement seule
        (comparer \qty{15}{\micro m} aux fortes bandes de l'eau
        au-dessous de \qty{8}{\micro m} et au-dessus de
        \qty{20}{\micro m})~?
\end{enumerate}

\medskip\noindent\textbf{Partie II --- L'échelle quantique de la flexion.}
\begin{enumerate}[resume]
  \item Calculer $\hbar\omega_2$ en meV, ainsi que l'échelle $E_n$.
  \item Quelle fraction des molécules occupe $n = 1$ à \qty{288}{K}
        (relativement à $n = 0$, facteur de Boltzmann)~? Et $n = 2$~?
  \item Quelle longueur d'onde de photon provoque $n = 0 \to 1$~?
        Pourquoi la \emph{même} longueur d'onde domine-t-elle
        l'émission~?
  \item Justifier, à partir de l'échelle harmonique, qu'une seule
        longueur d'onde serve toute l'échelle ($n \to n + 1$ pour tout
        $n$)~: quelle propriété de l'espacement des niveaux joue~?
  \item La molécule tourne aussi, ce qui ajoute une structure fine~:
        la signature à \qty{15}{\micro m} est en réalité une
        \emph{bande} large d'environ \qty{1}{\micro m}. Pourquoi la
        \emph{largeur} de bande importe-t-elle pour un gaz à effet de
        serre (penser à ce qui arrive une fois le centre de bande
        opaque)~?
  \item Un CO$_2$ vibrationnellement excité dans l'air a bien plus de
        chances de perdre son quantum par collision que par
        rayonnement (durée de vie radiative $\sim\qty{1}{s}$, temps
        entre collisions $\sim\qty{e-9}{s}$)~: où passe réellement
        l'énergie rayonnante absorbée~?
  \item Inversement, l'air à \qty{288}{K} entretient une population
        thermique dans $n = 1$ (question 8)~: que fait cette
        population, qui importe au bilan énergétique~?
\end{enumerate}

\medskip\noindent\textbf{Partie III --- Le grand livre radiatif de la planète.}
\begin{enumerate}[resume]
  \item Le Soleil rayonne comme un corps à \qty{5800}{K}~: calculer
        son pic de Wien. La flexion du CO$_2$ intercepte-t-elle
        beaucoup de lumière solaire~?
  \item Le sol rayonne comme un corps à \qty{288}{K}~: calculer son
        pic de Wien, et situer \qty{15}{\micro m} sur cette courbe
        thermique.
  \item Expliquer le mécanisme de l'effet de serre en quatre phrases~:
        le soleil entre, l'infrarouge thermique sort, l'interception à
        \qty{15}{\micro m}, la réémission vers le haut \emph{et} vers
        le bas.
  \item La part de la réémission qui s'échappe vers l'espace provient
        de couches hautes et froides ($\sim\qty{220}{K}$)~: pourquoi
        émettre depuis une couche plus froide réduit-il la puissance
        sortante de la planète à ces longueurs d'onde (rappeler
        que l'émission thermique croît avec $T$)~?
  \item Davantage de CO$_2$ pousse la couche émettrice plus haut et
        plus froid~: dites en une phrase pourquoi la surface doit alors
        se réchauffer pour rééquilibrer les comptes.
  \item Le centre de bande est déjà opaque~; l'effet du CO$_2$ ajouté
        agit dans les \emph{ailes} de la bande, d'où une croissance
        logarithmique du forçage avec la concentration~: relier cela à
        la question 11.
\end{enumerate}

\medskip\noindent\textbf{Partie IV --- Les isotopes~: l'empreinte
quantique.}
\begin{enumerate}[resume]
  \item La fréquence d'un oscillateur varie comme $\sqrt{k/\mu}$~:
        pour l'élongation antisymétrique, remplacer le $^{12}$C par du
        $^{13}$C change la masse effective~; la raie observée passe de
        \qty{2349}{cm^{-1}} à environ \qty{2283}{cm^{-1}}. Vérifier
        l'ordre de grandeur de ce décalage de $\approx 3\%$ à partir
        des masses.
  \item Des lasers accordés sur ces deux raies comptent séparément le
        $^{13}$CO$_2$ et le $^{12}$CO$_2$ dans un échantillon de gaz~:
        expliquer pourquoi l'échelle quantifiée rend possible une telle
        détection résolue en isotopes.
  \item Les plantes préfèrent l'isotope léger, si bien que le carbone
        fossile est pauvre en $^{13}$C~: qu'a fait le rapport
        isotopique mesuré du CO$_2$ atmosphérique à mesure que sa
        concentration montait --- et que prouve cela sur la
        \emph{source} du gaz ajouté~?
  \item La même physique à \qty{15}{\micro m} opère sur Vénus
        ($96\%$ de CO$_2$, \qty{90}{bar})~: qu'illustre sa surface à
        \qty{737}{K}~?
  \item Mars aussi respire du CO$_2$ presque pur, mais à \qty{6}{mbar},
        et elle est glaciale~: que démontre le trio
        Vénus--Terre--Mars sur la variable qui contrôle l'intensité de
        l'effet~?
  \item Résumer le résultat nommé~: un quantum de \qty{83}{meV} ---
        le $\hbar\omega$ d'une molécule triatomique qui se plie ---
        garé à \qty{15}{\micro m} sur le spectre thermique d'une
        planète à \qty{288}{K}, absorbé, thermalisé en quelques
        nanosecondes puis réémis depuis des altitudes froides, est le
        mécanisme par lequel $0.04\%$ de l'air fixe la température de
        la Terre.
\end{enumerate}
\end{problem}
